\documentclass{uic-beamer}
\usepackage{amsfonts,amsmath,oldgerm}

\newcommand{\testcolor}[1]{\colorbox{#1}{\textcolor{#1}{test}} \texttt{#1}}

\usefonttheme[onlymath]{serif}
\newcommand{\hrefcol}[2]{\textcolor{cyan}{\href{#1}{#2}}}

\title{UIC Presentation Style}
\subtitle{Using \LaTeX\ to prepare slides}
\author{\href{mailto:stephanie@math.uic.edu}{Stephanie Reyes}}
\date{Created on October 19, 2019}
\titlebackground{}

% changes \emph text to be bold and red

\begin{document}
\maketitle

\begin{frame}{Beamer for UIC slides}
\begin{itemize}
\item This is an \emph{intermediate} tutorial/template. 
\hrefcol{http://en.wikibooks.org/wiki/LaTeX/}{Here is one of many references available for learning LaTeX.}
\item Beamer has also a detailed
\hrefcol{http://www.ctan.org/tex-archive/macros/latex/contrib/beamer/doc/beameruserguide.pdf}{user manual}
\item \hrefcol{http://tug.ctan.org/macros/latex/contrib/beamer/doc/beameruserguide.pdf}{This reference} seems pretty decent. It has a lot of good notes on presentations in general.
\item \hrefcol{https://www.overleaf.com/learn/latex/Beamer}{Overleaf tutorials} are also \emph{really great} for learning \LaTeX.
\end{itemize}
\end{frame}

\begin{frame}{Beamer vs. PowerPoint}
\begin{itemize}
\item Math typesetting in \TeX\ is the best:
\begin{equation*}
\mathrm{i}\,\hslash\frac{\partial}{\partial t} \Psi(\mathbf{r},t) =
-\frac{\hslash^2}{2\,m}\nabla^2\Psi(\mathbf{r},t)
+ V(\mathbf{r})\Psi(\mathbf{r},t)
\end{equation*}

\end{itemize}
\end{frame}

%\begin{frame}[fragile]{Selecting the Class}
%\begin{block}{Minimum SINTEF Beamer Document}
%\verb|\documentclass{sintefbeamer}|\\
%\verb|\begin{document}|\\
%\verb|\begin{frame}{Hello, world!}|\\
%\verb|\end{frame}|\\
%\verb|\end{document}|\\
%\end{block}
%\end{frame}

%\begin{frame}[fragile]{Title page}
%To set a typical title page, you call some commands in the preamble:
%\begin{block}{The Commands for the Title Page}
%\begin{verbatim}
%\title{Sample Title}
%\subtitle{Sample subtitle}
%\author{First Author, Second Author}
%\date{Defaults to today's}
%\end{verbatim}
%\end{block}
%You can then write out the title page with %\verb|\maketitle|.

%You can set a different background image than the default %one with the
%\verb|\titlebackground| command, set before %\verb|\maketitle|.

%In the \texttt{backgrounds} folder, you can find a lot of %standard backgrounds
%for SINTEF presentation title pages.

%\end{frame}

\begin{frame}[fragile]{Writing a Simple Slide}
\framesubtitle{It's really easy!}
\begin{block}{Code for an Itemised List}<+->
\begin{verbatim}
\begin{frame}
  \begin{itemize}[<+->]
    \item ...
  \end{itemize}\end{verbatim}
\end{block}
\end{frame}

\begin{frame}[fragile]{Using Colours}
  \begin{itemize}
  \item Primary colour: \testcolor{uic-maintext};
  \item Contrast colours: \testcolor{uic-lblue}, \testcolor{uic-red}, 
                          \testcolor{uic-green}, \testcolor{uic-yellow};
  \item Additional colours: \testcolor{uic-grey}, \testcolor{uic-lgrey}.
  \end{itemize}
\end{frame}

\begin{frame}[fragile]{Adding images}
Adding images works like in normal \LaTeX:
\begin{columns}
\begin{column}{0.7\textwidth}
\begin{block}{Code for Adding Images}
\begin{verbatim}
\usepackage{graphicx}
% ...
\includegraphics
[width=\textwidth]
{uic-mscs-black}
\end{verbatim}
\end{block}
\end{column}
\begin{column}{0.3\textwidth}
\includegraphics
[width=\textwidth]{uic-mscs-black}\\
\end{column}
\end{columns}
\end{frame}

\begin{frame}[fragile]{Splitting in Columns}
\begin{columns}
\begin{column}{0.6\textwidth}
This is the first column
\end{column}
\begin{column}{0.3\textwidth}
And this the second
\end{column}
\end{columns}
\end{frame}

\begin{frame}[fragile]
\frametitle{Fonts}
  \begin{itemize}
  \item {\textrm{Use serif fonts only with high-definition projectors}}
  \item {\textsf{Use sans-serif fonts otherwise (or if you simply prefer them)}}
  \item This template uses the \emph{bookman} font. It is possible to use helvetica if you use a different compiler, but it may cause issues with other template settings.
  \end{itemize}
\end{frame}

\begin{frame}[fragile]{Look}
\begin{itemize}
\item This template can be adapted to a red theme using \hrefcol{https://marketing.uic.edu/marketing-toolbox/uic-fonts/}{uic fonts}
    and \hrefcol{https://marketing.uic.edu/marketing-toolbox/logos/}{uic logos} 
\item \emph{I fixed the light theme.} 
\item To insert a final slide, use \verb|\backmatter|.
\end{itemize}
\end{frame}


\begin{frame}[fragile]{Additional Notes}
\framesubtitle{incomplete}
\begin{itemize}
\item Don't load unnecessary packages. They may conflict with template and theme settings.
\end{itemize}
\end{frame}

\begin{frame}{Good Luck!}
\begin{itemize}
\item I know enough by now from working on this introduction to make my MS Thesis presentation.
\item If you have corrections or suggestions,
\hrefcol{mailto:stephanie@math.uic.edu}{send them to me!}\quad  
\end{itemize}
\end{frame}

\backmatter

\end{document}
